% Źródła: Eves s. 314--317 (§7.9); Wikipedia: Henri Poincaré, János Bolyai.

\section{Która geometria jest prawdziwa?} % lang-pl
\section{Quale geometria è quella vera?} % lang-it
\label{sekcja_fizyka}

Skoro istnieją trzy geometrie płaszczyzny, ktoś w końcu zapyta, która z nich opisuje przestrzeń, w której żyjemy. % lang-pl
W tej sekcji zobaczymy, dlaczego pomiar sumy kątów trójkąta nic nie rozstrzygnie, i poznamy świat wymyślony przez Poincarego, w którym mieszkańcy uwierzą, że żyją w przestrzeni Łobaczewskiego. % lang-pl
Poiché esistono tre geometrie del piano, prima o poi qualcuno chiederà quale di esse descriva lo spazio in cui viviamo. % lang-it
In questa sezione vedremo perché la misura della somma degli angoli di un triangolo non deciderà nulla, e conosceremo il mondo immaginato da Poincaré, i cui abitanti crederanno di vivere in uno spazio di Lobachevskij. % lang-it

Pytanie „która jest prawdziwa?'' znaczy: która dobrze opisuje przestrzeń fizyczną. % lang-pl
Wszystkie trzy geometrie pasują do naszego skromnego kawałka przestrzeni równie dobrze; jedyny rozstrzygający test -- pomiar sumy kątów dużego trójkąta -- dotąd nie wykazał odchylenia od $180^\circ$ większego niż błąd pomiaru \cite[s. 315]{eves1_1972}. % lang-pl
Odchylenie jest jednak proporcjonalne do pola trójkąta (zobacz sekcje \ref{sekcja_suma_katow} i \ref{sekcja_sferyczna}), a pole każdego dotąd zmierzonego trójkąta może być tak małe, że wszelka rozbieżność zniknie w dopuszczalnym błędzie. % lang-pl
Bolyai zauważy już w 1832 roku, że samym rozumowaniem matematycznym nie da się rozstrzygnąć, jaka jest geometria świata fizycznego (zobacz sekcję \ref{sekcja_gauss_bolyai}). % lang-pl
La domanda «quale è quella vera?» significa: quale descrive bene lo spazio fisico. % lang-it
Tutte e tre le geometrie si adattano altrettanto bene al nostro modesto pezzo di spazio; l'unico test decisivo -- la misura della somma degli angoli di un grande triangolo -- finora non ha mostrato uno scarto da $180^\circ$ superiore all'errore di misura \cite[p. 315]{eves1_1972}. % lang-it
Lo scarto è però proporzionale all'area del triangolo (si vedano le sezioni \ref{sekcja_suma_katow} e \ref{sekcja_sferyczna}), e l'area di ogni triangolo misurato finora può essere così piccola che qualsiasi discrepanza sparisce nell'errore ammesso. % lang-it
Bolyai noterà già nel 1832 che con il solo ragionamento matematico non si può decidere quale sia la geometria del mondo fisico (si veda la sezione \ref{sekcja_gauss_bolyai}). % lang-it

Poincaré pójdzie dalej: ponieważ każdy pomiar miesza założenia geometryczne z fizycznymi, obserwowany wynik da się wyjaśnić na wiele sposobów, wprowadzając odpowiednie poprawki do założeń o materii. % lang-pl
Wyobrazi więc sobie świat $\Sigma$, wypełniający wnętrze kuli o promieniu $R$ w przestrzeni euklidesowej, w którym obowiązują cztery prawa \cite[s. 315--316]{eves1_1972}: % lang-pl
Poincaré andrà oltre: poiché ogni misura mescola assunzioni geometriche e fisiche, il risultato osservato si può spiegare in molti modi, introducendo opportune correzioni nelle ipotesi sulla materia. % lang-it
Si immaginerà dunque un mondo $\Sigma$, che riempie l'interno di una sfera di raggio $R$ nello spazio euclideo, in cui valgono quattro leggi \cite[p. 315--316]{eves1_1972}: % lang-it
\begin{enumerate}
\item temperatura bezwzględna w odległości $r$ od środka wynosi $T = k(R^2 - r^2)$; % lang-pl
\item wymiary liniowe ciał są wprost proporcjonalne do temperatury w ich miejscu; % lang-pl
\item wszystkie ciała natychmiast przyjmują temperaturę swojego położenia; % lang-pl
\item światło biegnie po geodetykach świata $\Sigma$. % lang-pl
\item la temperatura assoluta a distanza $r$ dal centro è $T = k(R^2 - r^2)$; % lang-it
\item le dimensioni lineari dei corpi sono direttamente proporzionali alla temperatura del loro luogo; % lang-it
\item tutti i corpi assumono all'istante la temperatura della loro posizione; % lang-it
\item la luce si propaga lungo le geodetiche del mondo $\Sigma$. % lang-it
\end{enumerate}
Mieszkaniec $\Sigma$ nie wykryje tych praw ani odczuciem zmiany temperatury, ani termometrem, ani miarką, bo wszystko kurczy się razem z nim. % lang-pl
Nigdy nie dojdzie do brzegu, więc uzna swój świat za nieskończony, a geodetyki -- łuki okręgów prostopadłych do brzegu kuli -- wyda mu się prostymi, w których zachodzi aksjomat Łobaczewskiego (rysunek 7.9a u Evesa). % lang-pl
Zwykły kawałek przestrzeni euklidesowej będzie mu więc wyglądał na świat nieeuklidesowy \cite[s. 316--317]{eves1_1972}. % lang-pl
L'abitante di $\Sigma$ non scoprirà queste leggi né sentendo il cambiamento di temperatura, né con un termometro, né con un metro, perché tutto si restringe insieme a lui. % lang-it
Non arriverà mai al bordo, e quindi riterrà infinito il suo mondo, e le geodetiche -- archi di circonferenza ortogonali al bordo della sfera -- gli sembreranno rette, per le quali vale l'assioma di Lobachevskij (figura 7.9a in Eves). % lang-it
Un comune pezzo di spazio euclideo gli sembrerà dunque un mondo non euclideo \cite[p. 316--317]{eves1_1972}. % lang-it

Poincaré wyciągnie stąd wniosek, że pytanie o prawdziwą geometrię jest źle postawione: należy pytać, która jest \emph{najwygodniejsza}. % lang-pl
Do kreślenia, geodezji i budowy mostów będzie to geometria euklidesowa, bo najłatwiejsza. % lang-pl
Einstein użyje w ogólnej teorii względności jeszcze innej geometrii nieeuklidesowej -- tej, którą zapoczątkuje wykład Riemanna z 1854 roku (zobacz sekcję \ref{sekcja_riemann}) -- w której odstępstwa od geometrii euklidesowej zależą od rozkładu materii. % lang-pl
Badanie R.~K. Luneberga z 1947 roku wykaże z kolei, że przestrzeń widzenia najwygodniej opisać geometrią Łobaczewskiego \cite[s. 317]{eves1_1972}. % lang-pl
Poincaré ne trarrà la conclusione che la domanda sulla geometria vera sia mal posta: bisogna chiedersi quale sia la più \emph{comoda}. % lang-it
Per il disegno, la geodesia e la costruzione di ponti sarà la geometria euclidea, perché la più semplice. % lang-it
Einstein nella relatività generale userà ancora un'altra geometria non euclidea -- quella avviata dalla conferenza di Riemann del 1854 (si veda la sezione \ref{sekcja_riemann}) -- in cui gli scarti dalla geometria euclidea dipendono dalla distribuzione della materia. % lang-it
Lo studio di R.~K. Luneberg del 1947 mostrerà a sua volta che lo spazio visivo si descrive nel modo più comodo con la geometria di Lobachevskij \cite[p. 317]{eves1_1972}. % lang-it
\index[persons]{Einstein, Albert}%
\index[persons]{Luneberg, Rudolf K.}%
